% =====================================================================
%  05_bosquejo-general.tex — Requerimientos funcionales, arquitectura
%  general y base de datos.
% =====================================================================
\chapter{Bosquejo general}
\label{cap:bosquejo}

\section{Requerimientos funcionales}
\label{sec:rf}

Los requerimientos se organizan por módulo. Cada requerimiento principal
(\texttt{RF-NN}) agrupa requerimientos específicos (\texttt{RF-NN.M}). Los
valores numéricos corresponden a la configuración actual del
sistema~\cite{rep-tecnico}.

{\small
\begin{longtable}{@{}L{1.5cm}L{3.6cm}L{8.9cm}@{}}
  \caption{Requerimientos funcionales.}
  \label{tab:rf} \\
  \toprule
  \textbf{ID} & \textbf{Nombre} & \textbf{Descripción} \\
  \midrule
  \endfirsthead
  \caption[]{Requerimientos funcionales (continuación).} \\
  \toprule
  \textbf{ID} & \textbf{Nombre} & \textbf{Descripción} \\
  \midrule
  \endhead
  \textbf{RF-01} & \textbf{Adquisición de video} & El prototipo leerá de
    forma continua el video de la cámara de la puerta. \\
  RF-01.1 & Conexión a la cámara & El prototipo leerá el video por RTSP
    con decodificación por hardware y, si esta falla, usará FFmpeg. \\
  RF-01.2 & Reconexión automática & Si se pierde la señal, el prototipo
    intentará reconectarse cada 2~segundos sin detener los demás
    servicios. \\
  RF-01.3 & Modo prueba & El prototipo podrá usar archivos de video en
    lugar de la cámara, con controles para avanzar, retroceder o saltar
    entre videos. \\
  RF-01.4 & Segunda cámara & El prototipo permitirá cambiar a una segunda
    cámara opcional desde el tablero. \\
  \midrule
  \textbf{RF-02} & \textbf{Detección y seguimiento} & El prototipo
    detectará cada unidad presente en el cuadro y la seguirá mientras
    permanezca en escena. \\
  RF-02.1 & Detección de unidades & El prototipo detectará unidades con el
    modelo M1 y descartará las detecciones con confianza menor a 0.65. \\
  RF-02.2 & Seguimiento & El prototipo asignará a cada unidad un
    identificador de seguimiento y llevará por separado el número y el
    cruce de cada una, aunque haya varias a la vez. \\
  RF-02.3 & Límite de procesamiento & El prototipo limitará los cuadros
    por segundo de la detección a un valor configurable para no saturar
    la GPU. \\
  \midrule
  \textbf{RF-03} & \textbf{Lectura del número económico} & El prototipo
    obtendrá el número económico de cada unidad detectada. \\
  RF-03.1 & Localización de la zona & El prototipo localizará con el
    modelo M2 la zona del número (delantera, lateral o trasera) y la
    asignará a la unidad que la contiene. \\
  RF-03.2 & Lectura de dígitos & El prototipo leerá los dígitos de la zona
    con el modelo M3, previa mejora de contraste con CLAHE. \\
  RF-03.3 & Confirmación & El prototipo confirmará un número con 2 lecturas
    iguales si está en el padrón, o con 4 si no lo está. \\
  RF-03.4 & Número más largo & Si un número leído está contenido en otro
    más largo, el prototipo conservará el más largo. \\
  RF-03.5 & Topes de lectura & El prototipo hará como máximo 2 lecturas
    por unidad y 4 por cuadro, y releerá cada 2~segundos una unidad cuyo
    número ya esté confirmado. \\
  RF-03.6 & Recortes & El prototipo guardará recortes de las unidades para
    revisión y para ampliar los datos de entrenamiento. \\
  \midrule
  \textbf{RF-04} & \textbf{Determinación de cruces} & El prototipo decidirá
    si una unidad entra o sale del patio. \\
  RF-04.1 & Línea virtual & El prototipo comparará la posición de la base
    del recuadro de la unidad con una línea virtual configurable. \\
  RF-04.2 & Margen & El prototipo no cambiará el lado de una unidad que se
    encuentre a menos de 30~px de la línea. \\
  RF-04.3 & Maniobras & El prototipo descartará el movimiento de una unidad
    que regrese antes de 20~segundos. \\
  RF-04.4 & Repetición & El prototipo no registrará el mismo número dos
    veces en menos de 60~segundos. \\
  RF-04.5 & Registro único & Tras un cruce, el prototipo esperará a que la
    unidad termine de pasar (de 2 a 25~segundos) y registrará un solo
    evento con el mejor número y las mejores fotos. \\
  \midrule
  \textbf{RF-05} & \textbf{Registro de movimientos} & El prototipo guardará
    cada cruce y cada estancia en la base de datos. \\
  RF-05.1 & Evento de paso & El prototipo registrará la unidad, la
    dirección, la puerta, la zona leída, la confianza, la foto y la hora
    de cada cruce. \\
  RF-05.2 & Estancias & El prototipo abrirá una estancia con cada entrada,
    la cerrará con la salida correspondiente y calculará su duración. \\
  RF-05.3 & Padrón & El prototipo sincronizará las unidades válidas desde
    el archivo del padrón en cada arranque. \\
  RF-05.4 & Fuera del padrón & El prototipo registrará tal cual los
    números confirmados que no estén en el padrón, marcados como tales. \\
  RF-05.5 & Día operativo & El prototipo contará como un mismo día los
    movimientos entre las 3:00~a.m. y las 2:59~a.m. del día siguiente. \\
  \midrule
  \textbf{RF-06} & \textbf{Avisos y validación} & El prototipo notificará
    cada cruce a los responsables. \\
  RF-06.1 & Aviso por Telegram & El prototipo enviará al grupo de
    operación un aviso con la dirección, el tipo, la empresa y hasta 3
    fotos separadas al menos 1~segundo. \\
  RF-06.2 & Validación humana & El aviso incluirá botones para que un
    validador confirme o corrija la lectura. \\
  RF-06.3 & Validación automática & El prototipo podrá enviar la foto a
    Gemini para revisar la lectura antes de registrarla. \\
  RF-06.4 & Aviso por WhatsApp & El prototipo reenviará los mismos avisos
    a un número de WhatsApp configurado. \\
  RF-06.5 & Modo sin avisos & El prototipo permitirá desactivar los avisos
    para pruebas. \\
  \midrule
  \textbf{RF-07} & \textbf{Tablero web} & El prototipo ofrecerá un tablero
    web de consulta. \\
  RF-07.1 & Indicadores & El tablero mostrará indicadores, gráficas,
    unidades dentro y estancias. \\
  RF-07.2 & Video en vivo & El tablero mostrará el video de la cámara a
    10~cuadros por segundo y 960~px de ancho. \\
  RF-07.3 & Lista de cruces & El tablero listará los cruces con búsqueda
    por número o empresa, filtro por fechas y páginas de 10. \\
  RF-07.4 & Trayectoria & El tablero mostrará el recorrido de una unidad. \\
  RF-07.5 & Estado del sistema & El prototipo informará su estado
    (conexión, cuadros por segundo, contadores) en una ruta de salud. \\
  \midrule
  \textbf{RF-08} & \textbf{Grabación continua} & El prototipo conservará el
    video de la cámara como evidencia. \\
  RF-08.1 & Grabación por hora & El prototipo grabará el video sin
    recodificar en archivos de una hora y conservará las últimas 72~horas.
    \\
  RF-08.2 & Consulta de grabaciones & El tablero permitirá consultar y
    reproducir las grabaciones por hora. \\
  \bottomrule
\end{longtable}
}

\section{Arquitectura general del prototipo}

\subsection{Arquitectura en capas}

\saru{} se organiza en capas, cada una con una responsabilidad. Esta
separación permite probar cada parte por separado, cambiar un servicio
externo sin tocar la lógica de visión y que varias personas trabajen en
módulos distintos sin pisarse.

\begin{figure}[htbp]
  \centering
  \resizebox{\linewidth}{!}{%
  \begin{tikzpicture}[
      capa/.style={draw=azulCISA, fill=azulClaro, rounded corners=3pt,
                   minimum width=14cm, minimum height=1.25cm, align=center,
                   font=\small},
      capaAlt/.style={capa, draw=verdeCISA, fill=verdeFondo},
      capaGris/.style={capa, draw=grisTexto, fill=grisFondo},
    ]
    \node[capa] (pres) at (0,0) {\textbf{Presentación}\quad tablero,
      video en vivo, grabaciones\quad(\texttt{web/}, \texttt{templates/},
      \texttt{static/})};
    \node[capaAlt] (logica) at (0,-1.6) {\textbf{Lógica de visión y cruce}\quad
      cámara, M1, M2+M3, votación, línea virtual\quad(\texttt{logica/},
      \texttt{ml/})};
    \node[capa] (infra) at (0,-3.2) {\textbf{Servicios externos}\quad
      Telegram, Gemini, WhatsApp\quad(\texttt{avisos/}, \texttt{whatsapp/})};
    \node[capaGris] (datos) at (0,-4.8) {\textbf{Datos}\quad SQLite, padrón,
      fotos, recortes y grabaciones\quad(\texttt{metricas/},
      \texttt{scripts/grabar.sh})};
    \draw[flecha] (pres) -- (logica);
    \draw[flecha] (logica) -- (infra);
    \draw[flecha] (infra) -- (datos);
  \end{tikzpicture}%
  }
  \caption{Arquitectura general del prototipo.}
  \label{fig:capas}
\end{figure}

\begin{enumerate}
  \item \textbf{Presentación}: Flask sirve el tablero, el video en vivo,
        las fotos, las grabaciones y una API JSON. Las rutas solo leen
        datos; no contienen lógica de visión ni de almacenamiento.
  \item \textbf{Lógica de visión y cruce}: la cámara publica el cuadro más
        reciente; un hilo ejecuta M1 y el seguimiento, y otro ejecuta M2 y
        M3. La línea virtual decide los cruces y elige el mejor número.
  \item \textbf{Servicios externos}: envío de avisos a Telegram y WhatsApp
        y validación con Gemini. Todos se pueden desactivar sin afectar el
        registro.
  \item \textbf{Datos}: un solo módulo accede a la base SQLite; el padrón se
        lee de un archivo CSV; fotos, recortes y grabaciones se guardan en
        el disco de la Jetson.
\end{enumerate}

El sistema se despliega con Docker Compose en tres servicios:
\texttt{cisa-vision} (todas las capas anteriores), \texttt{cisa-grabacion}
(grabación continua) y \texttt{cisa-whatsapp} (puente de WhatsApp). La
grabación es independiente de la detección: una falla en una no detiene a
la otra~\cite{rep-tecnico}.

\section{Diagrama de base de datos}

El modelo relacional se organiza en dos ejes: los \textbf{catálogos}
(tipos de unidad, puertas y unidades) y los \textbf{movimientos} (eventos
de paso y estancias). La duración y el estatus de cada estancia no se
escriben: la base los calcula a partir de las horas de entrada y salida.

\begin{figure}[htbp]
  \centering
  \resizebox{\linewidth}{!}{%
  \begin{tikzpicture}[
      tabla/.style={draw=azulCISA, fill=white, rounded corners=2pt,
                    align=left, font=\scriptsize\ttfamily, inner sep=4pt},
      fk/.style={-{Stealth[length=2mm]}, thick, draw=azulCISA},
    ]
    \node[tabla] (tipo) at (0,0) {\begin{tabular}{@{}l@{}}
      \textbf{tipo\_unidad}\\\hline
      id\_tipo (PK)\\codigo\\descripcion\end{tabular}};
    \node[tabla] (unidad) at (4.2,0) {\begin{tabular}{@{}l@{}}
      \textbf{unidad}\\\hline
      id\_unidad (PK)\\no\_economico\\empresa\\id\_tipo (FK)\\en\_padron
      \end{tabular}};
    \node[tabla] (puerta) at (0,-3.6) {\begin{tabular}{@{}l@{}}
      \textbf{puerta}\\\hline
      id\_puerta (PK)\\nombre\end{tabular}};
    \node[tabla] (evento) at (4.2,-3.6) {\begin{tabular}{@{}l@{}}
      \textbf{evento\_paso}\\\hline
      id\_evento (PK)\\id\_unidad (FK)\\direccion\\id\_puerta (FK)\\zona\\
      confianza\\ruta\_captura\\hora\_paso\\hora\_registro\end{tabular}};
    \node[tabla] (estancia) at (9.0,-1.8) {\begin{tabular}{@{}l@{}}
      \textbf{estancia}\\\hline
      id\_estancia (PK)\\id\_unidad (FK)\\entrada\_id (FK)\\salida\_id (FK)\\
      hora\_entrada\\hora\_salida\\minutos (calculada)\\estatus (calculada)
      \end{tabular}};
    \draw[fk] (unidad) -- (tipo);
    \draw[fk] (evento) -- (unidad);
    \draw[fk] (evento) -- (puerta);
    \draw[fk] (estancia) -- (unidad);
    \draw[fk] (estancia) -- (evento);
  \end{tikzpicture}%
  }
  \caption{Modelo relacional de la base de datos. Las flechas van de la
           llave foránea a la tabla referida.}
  \label{fig:bd}
\end{figure}

Sobre estas tablas se definen vistas para las consultas del tablero:
unidades dentro en este momento, historial y promedio de estancias por
empresa, y números detectados fuera del padrón.
